\setcounter{section}{15}

  \zwischenueberschrift{Soal latihan}

\inputaufgabegibtloesung
{}
{

Misalkan $M$ sebuah manifold diferensiabel kompak dengan suatu bentuk volume positif yang kontinu $\omega$. Tunjukkan bahwa

\mavergleichskettedisp
{\vergleichskette
{ \int_M \omega
}
{ <} { \infty
}
{ } {
}
{ } {
}
{ } {
}
}
{}{}{}
berlaku.

}
{} {}

\inputaufgabe
{}
{

Tunjukkan bahwa ukuran volume yang diperkenalkan dalam
[[Mannigfaltigkeit/Abzählbar/Positive Volumenform/Zugehöriges Maß/Definition|Definisi 15.3]]
untuk suatu
\definitionsverweis {bentuk volume positif}{}{}
\definitionsverweis {kontinu}{}{}
pada sebuah
\definitionsverweis {manifold diferensiabel}{}{}
merupakan suatu
\definitionsverweis {ukuran}{}{}
$\sigma$-hingga.

}
{} {}

\inputaufgabe
{}
{

Misalkan

\mavergleichskettedisp
{\vergleichskette
{ \omega
}
{ =} { dx_1 \wedge \ldots \wedge dx_n
}
{ =} { e_1^* \wedge \ldots \wedge e_n^*
}
{ } {
}
{ } {
}
}
{}{}{}
merupakan bentuk volume standar pada $\R^n$. Tunjukkan bahwa untuk setiap
\definitionsverweis {himpunan bagian terukur}{}{}

\mathl{T \subseteq \R^n}{} berlaku kesamaan

\mavergleichskettedisp
{\vergleichskette
{
\int_{ T  }   \omega
}
{ =} { \int_{  T  }     \,  d  \lambda^n
}
{ =} { \lambda^n(T)
}
{ } {
}
{ } {
}
}
{}{}{.}

}
{} {}

\inputaufgabe
{}
{

Misalkan $M$ sebuah
\definitionsverweis {manifold diferensiabel}{}{}
dengan suatu
\definitionsverweis {bentuk volume positif}{}{}
$\omega$. Misalkan
\mathl{T \subseteq M}{}
\definitionsverweis {terukur}{}{}
dan
\mathl{N \subseteq M}{} suatu
\definitionsverweis {himpunan nol}{}{.}
Tunjukkan bahwa

\mavergleichskettedisp
{\vergleichskette
{
\int_{ T  }   \omega
}
{ =} {
\int_{ T \setminus N  }   \omega
}
{ } {
}
{ } {
}
{ } {
}
}
{}{}{}
berlaku.

}
{} {}

\inputaufgabe
{}
{

Misalkan $M$ sebuah
\definitionsverweis {manifold diferensiabel}{}{}
\definitionsverweis {berdimensi}{}{} $n$
dengan suatu
\definitionsverweis {basis topologi terhitung}{}{}
dan misalkan
\mathkor {} {\omega_1} {dan} {\omega_2} {}
merupakan
\definitionsverweis {bentuk volume positif}{}{}
pada $M$. Tunjukkan bahwa untuk setiap
\definitionsverweis {himpunan bagian terukur}{}{}

\mathl{T \subseteq M}{} dan
\mathl{a,b \in \R_+}{} berlaku hubungan

\mavergleichskettedisp
{\vergleichskette
{
\int_{ T  }  (a \omega_1 + b \omega_2)
}
{ =} { a
\int_{  T  }   \omega_1+ b
\int_{  T  }   \omega_2
}
{ } {
}
{ } {
}
{ } {
}
}
{}{}{.}

}
{} {}

\inputaufgabe
{}
{

Misalkan $M$ sebuah
\definitionsverweis {manifold diferensiabel}{}{}
dengan
\definitionsverweis {basis topologi terhitung}{}{.}
Jelaskan bagaimana mendefinisikan \anfuehrung{himpunan nol}{} pada $M$ dengan mengacu pada peta, tanpa diberikan suatu
\definitionsverweis {ukuran}{}{.}
Tunjukkan pula bahwa jika suatu
\definitionsverweis {bentuk volume positif}{}{}
diberikan, maka himpunan-himpunan nol tersebut juga merupakan
\definitionsverweis {himpunan nol}{}{}
dalam pengertian teori ukuran.

}
{} {}

\inputaufgabe
{}
{

Nyatakan lingkaran satuan
\mathl{S^1 \subset \R^2}{} sebagai serat suatu pemetaan
\maabbdisp {} {\R^2} {\R
} {}
sedemikian sehingga bentuk volume $\omega$ yang diberikan menurut
[[Differenzierbare reguläre Abbildung/R^n/Faser besitzt Volumenform über Gradienten/Fakt|Korolari 15.6]]
bersifat positif. Hitunglah
\mathl{\int_{S^1} \omega}{.}

}
{} {}

\inputaufgabe
{}
{

Nyatakan permukaan bola satuan
\mathl{S^2 \subset \R^3}{} sebagai serat suatu pemetaan
\maabbdisp {} {\R^3} {\R
} {}
sedemikian sehingga bentuk volume $\omega$ yang diberikan menurut
[[Differenzierbare reguläre Abbildung/R^n/Faser besitzt Volumenform über Gradienten/Fakt|Korolari 15.6]]
bersifat positif. Hitunglah
\mathl{\int_{S^2} \omega}{.}

}
{} {}

\inputaufgabe
{}
{

Misalkan
\mathkor {} {L} {dan} {M} {}
merupakan dua
\definitionsverweis {manifold diferensiabel}{}{}
dan
\maabbdisp {\varphi} { L } { M
} {}
suatu
\definitionsverweis {pemetaan diferensiabel}{}{.}
Misalkan

\mavergleichskette
{\vergleichskette
{ \omega
}
{ \in }{
{ \mathcal E }^{  1 } (  M   )
}
{  }{
}
{    }{
}
{   }{
}
}
{}{}{}
sebuah bentuk diferensial terukur dengan
\definitionsverweis {bentuk diferensial tarik balik}{}{}

\mavergleichskette
{\vergleichskette
{ \varphi^*\omega
}
{ \in }{
{ \mathcal E }^{  1 } (  L   )
}
{  }{
}
{    }{
}
{   }{
}
}
{}{}{}
dan misalkan
\maabbdisp {\gamma} { I } { L
} {}
suatu
\definitionsverweis {kurva yang terdiferensialkan secara kontinu}{}{}
\zusatzklammer {$I$ suatu
\definitionsverweis {interval riil}{}{}} {} {.}
Tunjukkan bahwa untuk
\definitionsverweis {integral lintasan}{}{}
berlaku kesamaan

\mavergleichskettedisp
{\vergleichskette
{ \int_\gamma \varphi^* \omega
}
{ =} { \int_{\varphi \circ \gamma} \omega
}
{ } {
}
{ } {
}
{ } {
}
}
{}{}{.}

}
{} {}

\inputaufgabe
{}
{

Diberikan
\maabbeledisp {\gamma} { [0,2\pi] } {\R^2
} { t } {( \cos  t , \sin  t )
} {,}
hitunglah
\definitionsverweis {integral lintasan}{}{}
sepanjang lintasan ini untuk
\definitionsverweis {bentuk diferensial}{}{}
berikut:

a)
\mathl{xdx +ydy}{,}

b)
\mathl{xdx -ydy}{,}

c)
\mathl{ydx +xdy}{,}

d)
\mathl{ydx -xdy}{.}

}
{} {}

\inputaufgabegibtloesung
{}
{

Hitunglah integral lintasan
\mathl{\int_\gamma \omega}{} untuk
\maabbeledisp {\gamma} { [-1,0] } {\R^3
} { t } {(-t^2,t^3-1,t+2)
} {,}
dengan bentuk diferensial berderajat $1$

\mavergleichskettedisp
{\vergleichskette
{ \omega
}
{ =} { x^3dx -yzdy +xz^2 dz
}
{ } {
}
{ } {
}
{ } {
}
}
{}{}{}
pada $\R^3$.

}
{} {}

\inputaufgabegibtloesung
{}
{

Kita meninjau dua pemetaan diferensiabel
\maabbeledisp {\gamma} { [1,2] } {\R^2
} { t } {(t,t^{-1})
} {,}
dan
\maabbeledisp {\varphi} {\R^2} {\R^3
} {(u,v)} {(u^2,uv,-u+v^2)
} {,}
serta bentuk diferensial

\mavergleichskettedisp
{\vergleichskette
{ \omega
}
{ =} { xdx-zdy+dz
}
{ } {
}
{ } {
}
{ } {
}
}
{}{}{}
pada $\R^3$.

a) Hitunglah bentuk diferensial tarik balik
\mathl{\varphi^* \omega}{} pada $\R^2$.

b) Hitunglah integral lintasan dari bentuk diferensial
\mathl{\varphi^* \omega}{} sepanjang lintasan $\gamma$.

c) Hitunglah
\zusatzklammer {tanpa mengacu pada b)} {} {}
integral lintasan dari bentuk diferensial
\mathl{\omega}{} sepanjang lintasan $\varphi \circ \gamma$.

}
{} {}

\inputaufgabegibtloesung
{}
{

Kita meninjau dua pemetaan diferensiabel
\maabbeledisp {\gamma} { [1,c] } {\R^2
} { t } {(t,t^{3})
} {,}
\zusatzklammer {dengan \mathlk{c \geq 1}{}} {} {}
dan
\maabbeledisp {\varphi} {\R_+ \times \R_+ } {\R^3
} {(u,v)} {(u^3,u^2+v^2,u^{-1}v^{-1} )
} {,}
serta bentuk diferensial

\mavergleichskettedisp
{\vergleichskette
{ \omega
}
{ =} { (x-y)dx-z^2dy+dz
}
{ } {
}
{ } {
}
{ } {
}
}
{}{}{}
pada $\R^3$.

a) Hitunglah bentuk diferensial tarik balik
\mathl{\varphi^* \omega}{} pada $\R_+ \times \R_+$.

b) Hitunglah integral lintasan dari bentuk diferensial
\mathl{\varphi^* \omega}{} sepanjang lintasan $\gamma$ sebagai fungsi $c$.

}
{} {}

  \zwischenueberschrift{Soal untuk dikumpulkan}

\inputaufgabe
{4}
{

Misalkan $M$ sebuah
\definitionsverweis {manifold diferensiabel}{}{}
\definitionsverweis {berdimensi}{}{} $n$
dengan
\definitionsverweis {basis topologi terhitung}{}{.}
Misalkan $\omega$ suatu
\definitionsverweis {bentuk volume positif}{}{} pada $M$ dan misalkan $\mu$ suatu
\definitionsverweis {ukuran}{}{} pada $M$ yang didefinisikan oleh bentuk volume ini. Tunjukkan bahwa setiap
\definitionsverweis {submanifold tertutup}{}{} berdimensi
\mathl{\leq n-1}{} merupakan suatu
\definitionsverweis {himpunan nol}{}{.}

}
{} {}

\inputaufgabe
{4}
{

Misalkan

\mavergleichskette
{\vergleichskette
{a,b,c,d,r,s
}
{ \geq }{ 1
}
{  }{
}
{    }{
}
{   }{
}
}
{}{}{}
merupakan
\definitionsverweis {bilangan asli}{}{.}
Kita meninjau
\definitionsverweis {kurva yang terdiferensialkan secara kontinu}{}{}
\maabbeledisp {} { [0,1] } {\R^2
} { t } {(t^r,t^s)
} {.}
Hitunglah
\definitionsverweis {integral lintasan}{}{}
sepanjang lintasan ini untuk
\definitionsverweis {bentuk diferensial}{}{}

\mavergleichskette
{\vergleichskette
{ \omega
}
{ = }{ x^ay^b dx +x^c y^d dy
}
{  }{
}
{    }{
}
{   }{
}
}
{}{}{.}

}
{} {}

\inputaufgabe
{5}
{

Diberikan
\maabbeledisp {\gamma} { [0,2\pi] } {\R^3
} { t } {( \cos  t , \sin  t, t )
} {,}
hitunglah
\definitionsverweis {integral lintasan}{}{}
sepanjang lintasan ini untuk
\definitionsverweis {bentuk diferensial}{}{}

\mavergleichskette
{\vergleichskette
{ \omega
}
{ = }{ (y-z^3) dx +x^2dy -xzdz
}
{  }{
}
{    }{
}
{   }{
}
}
{}{}{.}

}
{} {}

 [[Kategorie:Latexseite]]
